% generated from Ctrl 1/hidden_eigs.tex
\chapter{Power-sum lower bounds on hidden eigenvalues in positive realizations}\label{chap:Ctrl1}
\chaptermark{Ctrl1}
\begin{chapterabstract}
\noindent
A positive realization of a transfer function of order $n$ is a nonnegative matrix of some
size $N\ge n$ carrying the $n$ poles among its eigenvalues; the $s=N-n$ remaining
eigenvalues are the \emph{hidden} eigenvalues, whose characterization Benvenuti and Farina
list as an open problem. We give a lower bound on $s$ from the power sums of the poles: if
$p_k$ denotes the $k$-th power sum and $\rho$ the spectral radius, then
$s\ge\max_{k\ge1}\lceil[-p_k]_+/\rho^k\rceil$. At $k=1$ this reproduces the quantity
$\zeta$ of Benvenuti's lower bound; for $k\ge2$ it is, as far as we can determine, new, and
it is strictly stronger. On the family $\M_m(\theta)=\{1\}\cup m\cdot\{\pm i\theta\}$ the
published lower bound is vacuous ($N\ge n$) while ours gives $N\ge 2m(1+\theta^{2})$; when
$\theta^{2}>1-1/2m$ this meets Benvenuti's own upper bound $4m$ and therefore determines
the minimal size exactly, where previously only $n\le N\le 4m$ was known. The bound is
attained: on $\{1\}\cup m\cdot\{\pm i\}$ it gives $N\ge4m$, and a direct sum of $m$ copies
of the $4$-cycle permutation matrix realizes $N=4m$. We also record two limitations. The
method has a hard ceiling --- it can never prove more than $N\le2n$ --- so it is
complementary to, not competitive with, the digraph bounds that give unbounded ratios; and
the ratio $2$ itself is not new territory, being exactly what the Benvenuti--Farina family
of 1999 achieves. Finally we explain why no such bound can exist in the
Boyle--Handelman/Kim--Ormes--Roush setting: there one appends only zeros, and zeros change
no power sum. All claims are machine-verified in exact rational arithmetic.
\end{chapterabstract}

\section{Introduction}\label{Ctrl1:sec:intro}

Let $H(z)$ be a rational transfer function of order $n$ with pole multiset
$\Lambda=\{\lambda_1,\dots,\lambda_n\}$ and spectral radius
$\rho=\max_i|\lambda_i|$. A \emph{positive realization} of $H$ is a triple
$(A,b,c)$ with $A\in\R_{\ge0}^{N\times N}$, $b,c\ge0$ and
$H(z)=c(zI-A)^{-1}b$. Existence is settled \cite{Farina96,ADFB96,OMK84}, and a clean
criterion is available in terms of an invariant polyhedral cone \cite{BF04}. What is not
settled is \emph{size}. The poles of $H$ are a subset of the eigenvalues of any positive
realization, and the remaining $s=N-n$ eigenvalues --- the \emph{hidden eigenvalues} --- are
the reason $N$ may exceed $n$. Benvenuti and Farina put it as follows
\cite[\S IX]{BF04}: ``An important issue related to minimality of positive systems is the
study of `hidden' eigenvalues, i.e.\ of the eigenvalues which are not poles and that
possibly one has to add in order to obtain a minimal positive realization. A full
characterization of this property may lead to a deeper and valuable insight into the
problem.'' The minimality problem remains open; a recent statement of its status is
\cite{TS25}.

This note contributes one inequality. It is elementary, and its only merit is that it
appears not to have been written down beyond its first instance.

\paragraph{The bound.} Write $p_k=p_k(\Lambda)=\sum_{i}\lambda_i^{k}$ for the $k$-th power
sum of the poles and $[x]_+=\max(x,0)$. Theorem~\\ref{Ctrl1:thm:main} states
\begin{equation}\label{Ctrl1:eq:bound}
s\;\ge\;\Bigl\lceil\;\max_{k\ge1}\;\frac{[-p_k]_+}{\rho^{k}}\;\Bigr\rceil .
\end{equation}

\paragraph{What is already known.} The case $k=1$ of \\eqref{Ctrl1:eq:bound} is published.
Benvenuti \cite[Thm.~4.7]{Ben20ELA} proves $N\ge N_D+\sum_{|s|<1}m(s)+\zeta$ with
$\zeta=\lceil-m_{\mathcal R_1}-\sum_{|s|<1}s\,m(s)\rceil$ when that quantity is positive,
and the proof is exactly the $k=1$ argument: ``since the trace of a nonnegative matrix is
nonnegative \dots\ since each one of these additional eigenvalues has to be not greater
than $1$, then their number is not lesser than $\zeta$''
\cite[pp.~379--380]{Ben20ELA}. The companion paper \cite{Ben20SCL} transfers the same
quantity to the realization dimension. We have found no power sum beyond the first
anywhere in this literature: not in \cite{Ben20ELA,Ben20SCL}, not in the digraph bound of
Nagy and Matolcsi \cite{NM03}, not in \cite{BF04,TS25}. Section~\\ref{Ctrl1:sec:improve} shows
that $k=2$ already does strictly better.

\paragraph{What we do not claim.} Two disclaimers, stated here rather than buried.
\begin{enumerate}
\item[(a)] The ratio $N/n=2$ is \emph{not} new territory. The family
$H(z,m)=1/\bigl((z-1)(z^{2^{m}}+1)\bigr)$ of Benvenuti and Farina \cite{BF99} has
$n=2^{m}+1$ and minimal positive realization of dimension $N=2^{m+1}=2n-2$, so ratio $2$ is
exactly what they already achieve; and Nagy and Matolcsi \cite{NM03} obtain \emph{unbounded}
ratios, with $n=4$ fixed and $N\ge m+1$ for arbitrary $m$.
\item[(b)] Worse, ratio $2$ is a hard ceiling for the present method.
Since $p_k\ge-\sum_i|\lambda_i|^{k}\ge-n\rho^{k}$, the right-hand side of \\eqref{Ctrl1:eq:bound}
never exceeds $n$, so \\eqref{Ctrl1:eq:bound} can never prove more than $N\le 2n$
(Proposition~\\ref{Ctrl1:prop:ceiling}). The bound is therefore complementary to the combinatorial
bounds of \cite{NM03}, not a competitor to them.
\end{enumerate}
What we do claim is (i) the inequality for $k\ge2$, (ii) that it strictly improves the
published lower bound, and (iii) that on an explicit family it closes the published gap
completely.

\section{The bound}\label{Ctrl1:sec:bound}

\begin{theorem}\label{Ctrl1:thm:main}
Let $A\in\R_{\ge0}^{N\times N}$ and let $\Lambda$ be a sub-multiset of $\operatorname{spec}A$
of size $n$ containing an eigenvalue of modulus $\rho(A)$; write $\rho=\rho(A)$ and
$s=N-n$. Then for every $k\ge1$
\[
s\;\ge\;\frac{-p_k(\Lambda)}{\rho^{k}},
\qquad\text{hence}\qquad
s\;\ge\;\Bigl\lceil\max_{k\ge1}\frac{[-p_k(\Lambda)]_+}{\rho^{k}}\Bigr\rceil .
\]
\end{theorem}

\begin{proof}
Let $\M$ be the complementary multiset, $|\M|=s$, so that
$\operatorname{spec}A=\Lambda\uplus\M$. Every eigenvalue of $A$ has modulus at most
$\rho(A)=\rho$, so $\bigl|\sum_{\mu\in\M}\mu^{k}\bigr|\le\sum_{\mu\in\M}|\mu|^{k}\le
s\rho^{k}$. Since $A$ is nonnegative, $\Tr(A^{k})\ge0$: it is a sum of products of
nonnegative entries. Therefore
\[
0\;\le\;\Tr(A^{k})\;=\;p_k(\Lambda)+\sum_{\mu\in\M}\mu^{k}\;\le\;p_k(\Lambda)+s\rho^{k},
\]
which rearranges to the claim; $s$ is an integer, so the ceiling may be taken.
\end{proof}

The hypothesis that $\Lambda$ contains a dominant eigenvalue is what makes $\rho$ computable
from the poles alone; in the realization setting it holds because $\rho(H)$ is a pole
\cite[Lem.~7]{BF04}.

\begin{proposition}[$k=1$ is Benvenuti's $\zeta$]\label{Ctrl1:prop:k1}
At $k=1$, \\eqref{Ctrl1:eq:bound} reads $s\ge\lceil-p_1(\Lambda)/\rho\rceil$, which after
normalising $\rho=1$ and splitting $\Lambda$ into its dominant and non-dominant parts is
the quantity $\zeta$ of {\rm\cite[Thm.~4.7]{Ben20ELA}}.
\end{proposition}

\begin{proof}
Normalise $\rho=1$. The dominant part of $\Lambda$ contributes $m_{\mathcal R_1}$ to
$p_1$, since the $h$-th roots of unity sum to zero for $h>1$
\cite[proof of Thm.~4.7]{Ben20ELA}, and the non-dominant part contributes
$\sum_{|s|<1}s\,m(s)$. Hence $-p_1=-m_{\mathcal R_1}-\sum_{|s|<1}s\,m(s)$, and the ceiling
of its positive part is $\zeta$.
\end{proof}

\begin{proposition}[Ceiling of the method]\label{Ctrl1:prop:ceiling}
The right-hand side of \\eqref{Ctrl1:eq:bound} is at most $n$; consequently
Theorem~\\ref{Ctrl1:thm:main} can never certify $N>2n$.
\end{proposition}

\begin{proof}
$-p_k\le\sum_i|\lambda_i|^{k}\le n\rho^{k}$.
\end{proof}

\section{Strict improvement, and an exact determination}\label{Ctrl1:sec:improve}

Fix $0<\theta<1$ and $m\ge1$, and let
\[
\M_m(\theta)=\{1\}\;\cup\;m\ \text{copies of}\ \{\,i\theta,\,-i\theta\,\},
\qquad n=2m+1,\ \rho=1 .
\]
Only $1$ is dominant, so in the notation of \cite{Ben20ELA} we have $r=1$,
$N_D=1$ and $m_{\mathcal R_1}=1$.

\begin{theorem}\label{Ctrl1:thm:improve}
For $\M_m(\theta)$:
\begin{enumerate}
\item[(i)] the lower bound of {\rm\cite[Thm.~4.7]{Ben20ELA}} gives $N\ge 2m+1=n$, i.e.\ it is
vacuous, because $\zeta=0$;
\item[(ii)] Theorem~\\ref{Ctrl1:thm:main} at $k=2$ gives
$s\ge\lceil 2m\theta^{2}-1\rceil$, hence $N\ge 2m+1+\lceil2m\theta^{2}-1\rceil$, which
exceeds $n$ whenever $\theta^{2}>1/2m$;
\item[(iii)] the upper bound of {\rm\cite[Thm.~4.4]{Ben20ELA}} gives $N\le 4m$;
\item[(iv)] if $\theta^{2}>1-\dfrac{1}{2m}$ then the bounds in (ii) and (iii) coincide and
\[
N_{\min}=4m .
\]
\end{enumerate}
\end{theorem}

\begin{proof}
(i) $p_1(\M_m(\theta))=1$ and the non-dominant numbers cancel in conjugate pairs, so
$m_{\mathcal R_1}+\sum_{|s|<1}s\,m(s)=1\ge0$ and $\zeta=0$; the bound reads
$N\ge N_D+2m+0=2m+1$. (ii) $p_2=1+m\bigl((i\theta)^{2}+(-i\theta)^{2}\bigr)=1-2m\theta^{2}$,
so $[-p_2]_+/\rho^{2}=2m\theta^{2}-1$ when positive. (iii) is \cite[Thm.~4.4]{Ben20ELA}
with $\kappa(i\theta)=4$, valid since $i\theta\in\Pi_4\setminus\Pi_3$ for
$\theta>1/\sqrt3$. (iv) The bound in (ii) equals $4m$ as soon as
$\lceil2m\theta^{2}-1\rceil\ge2m-1$, i.e.\ $2m\theta^{2}-1>2m-2$, i.e.\
$\theta^{2}>1-1/2m$.
\end{proof}

Table~\\ref{Ctrl1:tab:improve} shows the three quantities. The criterion $\theta^{2}>1-1/2m$ of
Theorem~\\ref{Ctrl1:thm:improve}(iv) predicts exactly which rows close: it holds for
$\theta=9/10$ only at $m=1$ ($0.81>0.5$) and fails at $m=4$ ($0.81<0.875$), where the bound
indeed falls one short of the upper bound; it holds throughout for $\theta=99/100$ from
$m=2$ on and for $\theta=999/1000$, and there the gap closes and $N$ is determined.

\begin{table}[ht]
\centering\small
\begin{tabular}{rrrrrr}
\toprule
$\theta$ & $m$ & $n$ & \cite[Thm.~4.7]{Ben20ELA} & Theorem~\\ref{Ctrl1:thm:main} & \cite[Thm.~4.4]{Ben20ELA}\\
 & & & lower & lower & upper\\
\midrule
$9/10$      &  1 &  3 &  3 &  \textbf{4} &  4\\
$9/10$      &  4 &  9 &  9 & 15 & 16\\
$9/10$      & 16 & 33 & 33 & 58 & 64\\
\midrule
$99/100$    &  2 &  5 &  5 &  \textbf{8} &  8\\
$99/100$    &  4 &  9 &  9 & \textbf{16} & 16\\
$99/100$    &  8 & 17 & 17 & \textbf{32} & 32\\
$99/100$    & 16 & 33 & 33 & \textbf{64} & 64\\
\midrule
$999/1000$  &  4 &  9 &  9 & \textbf{16} & 16\\
$999/1000$  & 16 & 33 & 33 & \textbf{64} & 64\\
\bottomrule
\end{tabular}
\caption{Lower and upper bounds on the minimal size $N$ for $\M_m(\theta)$. Bold: the new
lower bound meets the published upper bound, so $N$ is determined exactly. All entries
exact.}
\label{Ctrl1:tab:improve}
\end{table}

\section{Sharpness}\label{Ctrl1:sec:sharp}

Let $\M_m=\{1\}\cup m\cdot\{i,-i\}$, all of modulus one, $n=2m+1$.

\begin{theorem}\label{Ctrl1:thm:sharp}
$p_2(\M_m)=1-2m$, so Theorem~\\ref{Ctrl1:thm:main} gives $s\ge2m-1$ and $N\ge4m$. The bound is
attained: if $C_4$ is the $4$-cycle permutation matrix, with characteristic polynomial
$x^{4}-1$ and spectrum $\{1,i,-1,-i\}$, then $C_4^{\oplus m}$ is nonnegative of size $4m$
and contains $\M_m$ in its spectrum. Hence $N_{\min}=4m=2n-2$.
\end{theorem}

\begin{proof}
The power sum is immediate. For the construction, $\operatorname{spec}C_4^{\oplus m}$ is
$m$ copies of $\{1,i,-1,-i\}$, which contains one $1$ and $m$ copies of each of $\pm i$.
\end{proof}

\begin{remark}\label{Ctrl1:rem:caveat}
Two caveats. First, for $m\ge2$ the multiset $\M_m$ violates the standing hypothesis of
\cite[Thms.~3.2, 4.1]{Ben20ELA} that the spectral radius occur with dominant multiplicity:
here $m(1)=1$ while the dominant multiplicity is $m$. The formula of
\cite[Thm.~4.1]{Ben20ELA} nevertheless evaluates to $N_D=4m$ on $\M_m$, so the conclusion
$N_{\min}=4m$ is also reachable by that route, by an argument through roots of unity rather
than through traces; we note this so that Theorem~\\ref{Ctrl1:thm:sharp} is not mistaken for a new
determination. It is a sharpness check on \\eqref{Ctrl1:eq:bound}, not a new value of $N$. Second,
$4m=2n-2$ is the Benvenuti--Farina ratio of \cite{BF99}, as recalled in
Section~\\ref{Ctrl1:sec:intro}.
\end{remark}

\begin{remark}[Tutorial Example 8]\label{Ctrl1:rem:ex8}
For the poles $\{1,\pm0.9i\}$ of \cite[Ex.~8]{BF04}, $p_2=-31/50$ and \\eqref{Ctrl1:eq:bound}
gives $s\ge1$, i.e.\ $N\ge4$. This is attained by the nonnegative circulant with first row
$(0,\tfrac{19}{20},0,\tfrac1{20})$, whose characteristic polynomial is
$(x-1)(x+1)(x^{2}+\tfrac{81}{100})$: the hidden eigenvalue is $-1$. The example in
\cite{BF04} is $4$-dimensional, so the bound is tight there as well.
\end{remark}

\section{Why there is no analogue for the Boyle--Handelman problem}\label{Ctrl1:sec:bh}

It is worth saying why the same traces behave completely differently one category over.
In the inverse eigenvalue problem solved by Boyle and Handelman over $\R$ \cite{BH91} and by
Kim, Ormes and Roush over $\Z$ \cite{KOR00}, one asks whether $\Lambda$ is the \emph{nonzero}
spectrum of a nonnegative matrix; the matrix may be enlarged, but only by appending
\emph{zeros}. Zeros contribute nothing to any power sum, so the trace conditions
$\operatorname{tr}_k\ge0$ must already hold for $\Lambda$ itself: they are a genuine
obstruction, not a counting device, and no lower bound of the shape \\eqref{Ctrl1:eq:bound} on the
size of the dilation can exist. Equivalently: in that category the analogue of
Theorem~\\ref{Ctrl1:thm:main} is the statement $p_k(\Lambda)\ge0$, with no $s$ in it.

This is the exact point at which results from symbolic dynamics fail to transfer to the
positive realization problem. A negative $p_2$ closes the fixed-nonzero-spectrum problem
outright --- shift equivalence preserves the nonzero spectrum with multiplicities
\cite[\S7.4]{LM95}, so no dilation of any size helps, an argument used in this form in
\cite[II]{ANCFT} --- while in the present setting the same negative $p_2$ merely counts
hidden eigenvalues, and indeed \cite[Ex.~8]{BF04} exhibits a transfer function with
$p_2=-0.62$ and a $4$-dimensional positive realization. The two categories differ exactly by
the freedom to add nonzero eigenvalues.

\section{Verification}\label{Ctrl1:sec:battery}

The script \texttt{hidden\_eigs.py} accompanying this note runs the checks of
Table~\\ref{Ctrl1:tab:battery} in $23$\,s, all passing; its output is archived as
\texttt{hidden\_eigs\_output.txt}. Arithmetic is exact rational throughout; the only
floating-point quantity anywhere is the spectral radius in check (T), and that row is
labelled accordingly.

\begin{table}[ht]
\centering\small
\begin{tabular}{clc}
\toprule
key & check & mode\\
\midrule
(T) & Theorem~\\ref{Ctrl1:thm:main} holds on $120$ random nonnegative integer matrices, $\Lambda$ a rational factor & exact $+$ numeric $\rho$\\
(K) & at $k=1$ the bound equals $\zeta$ of \cite[Thm.~4.7]{Ben20ELA} & exact\\
(I) & $k=2$ strictly exceeds \cite[Thm.~4.7]{Ben20ELA} in $15/15$ instances (Table~\\ref{Ctrl1:tab:improve}) & exact\\
(S) & Theorem~\\ref{Ctrl1:thm:sharp}: bound $=4m$ attained by $C_4^{\oplus m}$, $m\le32$ & exact\\
(E) & Remark~\\ref{Ctrl1:rem:ex8}: bound $=1$, attained by an explicit rational circulant & exact\\
(C) & Proposition~\\ref{Ctrl1:prop:ceiling}: the bound never exceeds $n$, on $300$ random pole sets & exact\\
(Z) & appending zeros leaves every power sum unchanged (Section~\\ref{Ctrl1:sec:bh}) & exact\\
\bottomrule
\end{tabular}
\caption{The verification battery (\texttt{hidden\_eigs.py}), $7/7$ passing.}
\label{Ctrl1:tab:battery}
\end{table}

In check (T) the spectral radius is replaced by a rational \emph{lower} bound, which makes
the tested inequality strictly stronger than the theorem; no violation occurred. The bound
was attained in none of those $120$ random instances, which is the expected behaviour: like
the $k=1$ bound it is weak on generic spectra and bites only when the poles are strongly
oscillatory.

\section{Scope}\label{Ctrl1:sec:scope}

\emph{Established here.} Theorem~\\ref{Ctrl1:thm:main} for all $k$; its coincidence with
\cite[Thm.~4.7]{Ben20ELA} at $k=1$; strict improvement and an exact determination of
$N_{\min}$ on $\M_m(\theta)$ for $\theta^{2}>1-1/2m$; sharpness on $\M_m$ and on
\cite[Ex.~8]{BF04}; the ceiling $N\le2n$; the structural reason the bound has no analogue
in \cite{BH91,KOR00}.

\emph{Open.} (a) Whether power sums can be combined with the cone criterion \cite{BF04} or
with the digraph bound \cite{NM03} to pass the ratio-$2$ ceiling. (b) Whether the
Johnson--Loewy--London inequalities \cite{LoewyLondon78}, which mix power sums of different
orders, give a bound not implied by \\eqref{Ctrl1:eq:bound}; we have not investigated this and it
is the obvious next step. (c) Whether the bound is ever tight outside the families of
Sections~\\ref{Ctrl1:sec:improve}--\\ref{Ctrl1:sec:sharp}; it was not tight in any of our $120$ random
instances.

\emph{Not claimed.} Ratio $2$, which is \cite{BF99}; unbounded ratios, which are
\cite{NM03}; and the $k=1$ case, which is \cite{Ben20ELA}.

\section*{Provenance of this chapter}

Unrefereed preprint. Theorem~\\ref{Ctrl1:thm:main} is elementary and its proof is complete above.
The bounds of \cite{Ben20ELA} quoted for comparison in Table~\\ref{Ctrl1:tab:improve} are
evaluated by us from the published statements of Theorems 4.4 and 4.7 of that paper and are
reproduced in the script; we have not reproved them. The symbolic-dynamics facts in
Section~\\ref{Ctrl1:sec:bh} are quoted from \cite{LM95,BH91,KOR00}. We searched the literature for
a power-sum bound with $k\ge2$ and found none; a search cannot establish a negative, and we
state this as the limit of our knowledge.
